\section{Explicit choices for every exponent above three}\label{sec:parameters}

Fix a rational \(\nu>3\). Define rational numbers
\begin{align*}
 e&=\frac{\nu-3}{2(2\nu-3)},& s&=1-e,& v&=s^3,\\
 D&=\nu s^2-1-(\nu-1)s^3,&
 B&=s-\frac{D}{4\nu s},&
 C&=\frac12\left(\frac{B^2}{v}+\frac1B\right).
\end{align*}
Here \(0<s<1\), and
\begin{align*}
 D&=(1-s)((\nu-1)s^2-s-1)>0,\\
 (\nu-1)s^2-s-1&=\frac{\nu-3}{2}+(\nu-1)e^2>0.
\end{align*}
Moreover \(0<D<\nu s^2\), so \(0<B<s\), and
\begin{equation}\label{eq:G}
 G:=\nu B^2-1-(\nu-1)v
                   =\frac D2+\frac{D^2}{16\nu s^2}>0.
\end{equation}
It follows that
\begin{equation}\label{eq:BC}
 B^2>v,\quad B^3<v,\quad 1<B^2/v<C<1/B,
 \quad B^2<Cv<1,\quad CB<1.
\end{equation}

Choose a rational \(F>2/v\) large enough that \(\nu/F<G/8\).
For even integers \(m\) tending to infinity, set
\begin{equation}\label{eq:scaling}
 J=\lfloor C^{m/2}\rfloor,\quad K=J^2,\quad
 w_0=B^{-m},\quad v_0=2Jv^{m/2}B^{-m},\quad
 \sigma=\frac1{100(m+2)}.
\end{equation}
All these weights are rational and \(J\geq1\). Exactly,
\(K(w_0/v_0)^2v^m=1/4\). In fact the inflated volume is less than
\(25/97<1\), since
\((1+3/[100(m+2)])^{m+2}<1/(1-3/100)\).
Conditions \eqref{eq:shrink} and \eqref{eq:basegap} hold eventually:
by \eqref{eq:BC},
\begin{align*}
 K/w_0&\sim(CB)^m\longrightarrow0,\\
 v_0&\sim2(\sqrt{Cv}/B)^m\longrightarrow\infty,\\
 m/v_0&\longrightarrow0,\\
 w_0/v_0&\sim\tfrac12(Cv)^{-m/2}\longrightarrow\infty
             \quad\hbox{exponentially}.
\end{align*}
The last quantity eventually exceeds the quadratic expression
\(2(m+3)/\sigma\).

The value in \eqref{eq:lambda} simplifies exactly to
\[
 \lambda_*^{m+2}=\frac{(\nu B^2)^m}{64},
 \qquad \lambda_*\longrightarrow\nu B^2.
\]
Thus \(g_m\to G>0\). Choose and fix an even \(m\) large enough
for the geometric conditions and
\begin{equation}\label{eq:mchoice}
 g_m>G/2,\qquad
 E_0:=\frac{\Lambda Fm+2\log2}{v_0}+\frac{100K}{w_0}<G/8.
\end{equation}

Only now choose \(m\) fractions from the hypothetical unbounded
sequence \eqref{eq:exceptional}. Successively large denominators ensure
\(1<r_i<2\), \eqref{eq:qcondition}, all separation inequalities,
and
\begin{equation}\label{eq:wchoice}
 E_w:=\Lambda\sum_i\frac1{w_i}
        +\frac{v+\nu+\log12+\log(600K)}{w_*}<G/8.
\end{equation}
For the last condition it suffices first to require every \(w_i\)
to exceed one fixed large constant, and then enforce successive
separation. Unbounded denominators give unbounded
\(\lceil\log q\rceil\), even along the exceptional sequence.
No next fraction is needed when choosing a previous threshold.

Now all data are fixed. We have
\[
            E=\nu/F+E_0+E_w<3G/8<g_m.
\]
The interpolation proposition supplies the nonzero minors for all large
divisible \(H\). Taking \(H\) to infinity in their lower and upper
bounds gives \eqref{eq:contradiction}, contrary to \(g_m>E\). This excludes \eqref{eq:exceptional} for
each rational \(\nu>3\).

Given \(\varepsilon>0\), choose rational \(3<\nu<3+\varepsilon\).
The set \(\mathcal E\) of reduced fractions satisfying
\eqref{eq:exceptional} is finite: its denominators are bounded,
and the same inequality bounds its numerators.
For \(p/q\notin\mathcal E\), reduce the fraction to \(p'/q'\).
Then \(q'\leq q\), and
\[
 |\theta-p/q|=|\theta-p'/q'|>(q')^{-\nu}
                                  \geq q^{-\nu}\geq q^{-3-\varepsilon}.
\]
If \(\mathcal E\) is nonempty, then for its finitely many rational numbers,
irrationality of \(\theta\) gives a strictly positive minimum
distance from \(\theta\). Since \(q^{-3-\varepsilon}\) tends to
zero, increasing the cutoff covers every representation of these
fractions as well. If it is empty, no additional cutoff is needed.
This proves \eqref{eq:main} without a
coprimality restriction. Together with Dirichlet's theorem it
proves Theorem~\ref{thm:main}.
\hfill\(\square\)
